% ============================================================
% f2015ahw_q07_flowchart.tex
% 2015 安徽卷（文科）第 7 题 · 程序框图 —— 清单 F01 / §1
%
% 来源：.work/phase3-b12/figure-inventory.md §1（源 PDF 内容流逐 path 实测真页 pt）
%       源件 .work/phase3-b12/source/2015安徽文.pdf p1
%
% 坐标换算：画布 [x=1pt, y=1pt]（单位 TeX pt）。局部坐标 = 真页 pt − 原点
%   O = (125.387, 120.228)（清单 §1 几何 bbox 左下角）。真页 pt → 局部 pt 为
%   1:1 平移、不缩放——清单 §0.3 的 0.999622 是 pdftocairo 渲染件的 SVG 观测口径换算，
%   只用于与渲染件对读，不入绘图坐标。渲染层面 1 TeX pt = 0.996264 PDF pt（72/72.27），
%   整图统一 0.37%，属排版口径而非坐标换算误差。
%
% 线宽（清单 §0.4 本卷实测分档）：图形轮廓与连线 0.398pt；框内分式横线 0.42pt
% 虚线：无（本图全部实线）
% 箭头：arrows.meta Triangle 长 3.985pt、宽 2.988pt（清单 §6.3 全卷 31 枚同规格）；
%   箭尖落在清单尖端坐标；线端由 TikZ 按箭长回缩（源内线端再多深入 0.398pt = 一个
%   线宽，被填充三角完全覆盖，不可见）。
% 标签：字号 absolute 10.5pt（汉字 STSong 10.5（本图以 ctex 中文字体呈现）、数学 CMMI/CMR 10.5）；一律 anchor=base 落在清单实测基线上——
%   清单给的是字框中心：汉字字框 = 全角 em 方框（基线 = 中心 − 0.36em），拉丁/数学字框
%   = 字体 ascent/descent 方框（基线 = 中心 − 0.25em）；两式均与清单另给的基线自洽。
%
% 照源复刻的存疑项：
%    - 开始框高 16.838 / 结束框高 17.016，同宽 27.999 而高差 0.178pt —— 源内不一致，照源（B-S2）
%    - 回注折线 x 在 296.483/296.484 间有位级抖动、#22 段起讫 x 差 0.001pt（三位小数舍入），照源保留
%    - 箭头尖端与目标框之间一致留 0.199–0.282pt 间隙（源画法，非贴边），照源
%    - 分式标签 a=1+1/(1+a) 按源 600 dpi 墨迹实测逐符号落位（清单只给词块中心；尾随
%      二元算符在 TeX 数学列表中会失掉前导 glue，整块单节点会比源窄 2.3pt，故拆为逐
%      符号节点）。菱形条件按清单 §1.2 的 5 个词块中心逐块落位，块内间距由 TeX 自带；
%      本图字形（Latin Modern）与源（CM）度量略有出入，该行整体较源窄约 2.1pt，
%      块中心残差 ≤0.4pt。
% ============================================================
\begin{tikzpicture}[x=1pt, y=1pt,
  every node/.style={inner sep=0, font=\fontsize{10.5pt}{12.6pt}\selectfont},
  line cap=butt, line join=miter,
  arr/.style={-{Triangle[length=3.985pt, width=2.988pt]}},
  flow/.style={line width=0.398pt}]

  % ---- 节点（清单 §1.1 表）----
  \draw[flow, rounded corners=3.985pt] (71.041,166.194) rectangle (99.040,183.032);
  \draw[flow] (53.156,135.830) rectangle (116.924,151.623);
  \draw[flow] (158.933,102.764) -- (85.040,121.175) -- (11.146,102.764) -- (85.040,84.353) -- cycle;
  \draw[flow] (49.711,40.751) rectangle (120.369,69.698);
  % 框内分式横线（线宽 0.42pt）
  \draw[line width=0.42pt] (92.040,54.953) -- (115.673,54.953);
  \draw[flow] (59.196,11.538) rectangle (110.884,26.179);
  \draw[flow] (146.059,46.795) -- (152.195,63.654) -- (196.129,63.654) -- (189.993,46.795) -- cycle;
  \draw[flow, rounded corners=3.985pt] (157.098,15.206) rectangle (185.097,32.222);

  % ---- 连线与箭头（清单 §1.1；箭尖 = 清单尖端坐标）----
  \draw[flow, arr] (85.040,165.995) -- (85.040,151.821);
  \draw[flow, arr] (85.040,135.630) -- (85.040,121.457);
  \draw[flow, arr] (85.040,84.071) -- (85.040,69.897);
  \draw[flow, arr] (85.040,40.552) -- (85.040,26.378);
  \draw[flow, arr] (171.097,46.595) -- (171.096,32.421);
  % 「否」支路：菱形右顶点 → 右行 → 下行至平行四边形
  \draw[flow, arr] (159.215,102.764) -- (171.096,102.764) -- (171.096,63.855);
  % 回注环流：n=n+1 ↓ → 左 → 上 → 右，箭尖落在菱形上方竖直连线上
  \draw[flow, arr] (85.040,11.339) -- (85.040,0.000) -- (0.000,0.000) -- (0.000,128.544) -- (85.040,128.544);

  % ---- 标签（清单 §1.2；anchor=base 于实测基线）----
  \node[anchor=base] at (85.043,171.412) {开始};
  % a=1, n=1：按清单 §1.2 的 6 个词块中心逐块落位（同一基线 y=261.58）
  \node[anchor=base] at (59.433,141.352) {$a$};
  \node[anchor=base] at (69.203,141.352) {$=$};
  \node[anchor=base] at (80.283,141.352) {$1,$};
  \node[anchor=base] at (91.023,141.352) {$n$};
  \node[anchor=base] at (101.173,141.352) {$=$};
  \node[anchor=base] at (110.793,141.352) {$1$};
  % 菱形条件：按清单 §1.2 的 5 个词块中心逐块落位（同一基线 y=220.37）
  \node[anchor=base] at (43.773,100.142) {$|a$};
  \node[anchor=base] at (54.423,100.142) {$-$};
  \node[anchor=base] at (74.253,100.142) {$1.414|$};
  \node[anchor=base] at (94.673,100.142) {$\geqslant$};
  \node[anchor=base] at (116.103,100.142) {$0.005?$};
  \node[anchor=base] at (159.213,107.952) {否};
  \node[anchor=base] at (93.993,73.892) {是};
  % 分式标签 a=1+1/(1+a)：按源 600 dpi 墨迹实测逐符号落位（尾随二元算符
  %   在 TeX 中失 glue，整块单节点会少 2.3pt 空隙；故主行与分母均逐符号定位）
  \node[anchor=base] at (55.953,52.332) {$a$};
  \node[anchor=base] at (65.673,52.332) {$=$};
  \node[anchor=base] at (75.333,52.332) {$1$};
  \node[anchor=base] at (84.333,52.332) {$+$};
  \node[anchor=base] at (103.853,59.432) {$1$};
  \node[anchor=base] at (94.653,45.122) {$1$};
  \node[anchor=base] at (103.593,45.122) {$+$};
  \node[anchor=base] at (112.903,45.122) {$a$};
  \node[anchor=base] at (171.093,52.062) {输出$n$};
  \node[anchor=base] at (85.038,15.912) {$n=n+1$};
  \node[anchor=base] at (171.093,20.542) {结束};
\end{tikzpicture}
